%!TEX root = ../main.tex

\section{Training Framework Details}
\label{sec:appendix_training_details}

In this section, we provide additional details on the Progressive Agentic Training framework described in Section~\ref{sec:agentic_training_framework}, including the cold-start curriculum stages, environment token masking, and hybrid reward design.

\subsection{Cold-Start Curriculum Stage 1: Operator Syntax Learning}
\label{subsec:appendix_stage1}

As described in Section~\ref{subsec:cold_start}, Stage 1 focuses on the syntax and semantics of individual operators. We provide additional details on the data construction process.

\stitle{Data Construction.}
Given source tables and a ground-truth operator pipeline $\mathcal{P}_{\text{total}} = (o^{(1)}, \dots, o^{(N)})$, we first execute these operators sequentially to collect a sequence of intermediate table sets $\{\mathcal{T}_{0}, \dots, \mathcal{T}_{N}\}$. Then, we sample a contiguous sub-sequence of operators $\mathcal{P}_{\text{sub}} = (o^{(i)}, \dots, o^{(j)})$ with the corresponding input intermediate table set $\mathcal{T}_{\text{in}}=\mathcal{T}_{i-1}$ and output intermediate table set $\mathcal{T}_{\text{out}}=\mathcal{T}_{j}$.

\stitle{Sub-sequence Sampling Strategy.}
To ensure diverse and balanced training data, we employ the following sampling strategy:
\begin{itemize}[leftmargin=*]
    \item \textbf{Variable length.} We sample sub-sequences of varying lengths (1 to 5 operators) to expose the model to both single-operator and multi-operator transformations.
    \item \textbf{Random starting point.} For each ground-truth pipeline, we randomly select the starting position $i$ and then sample the ending position $j$ based on the desired sub-sequence length.
    \item \textbf{Coverage.} We ensure that all operator types appear in the training data with sufficient frequency by oversampling rare operator types.
\end{itemize}

Through this process, we construct the operator syntax learning dataset $\mathcal{D}_{\text{op}}$ consisting of training triples $(\mathcal{P}_{\text{sub}}, \mathcal{T}_{\text{in}}, \mathcal{T}_{\text{out}})$, and optimize the loss $\mathcal{L}_{\text{op}}$ defined in Eq.~(\ref{eq:sft_op_learn}).

\subsection{Cold-Start Curriculum Stage 2: Reasoning Procedure Learning}
\label{subsec:appendix_stage2}

Stage 2 aligns the model with the tree-based agentic reasoning mechanism. We provide additional details on the teacher model distillation process.

\stitle{Teacher Model Distillation.}
We employ five open-source LLMs as teacher models to generate diverse reasoning trajectories. For each ADP task defined by source tables $\mathcal{S}$ and a target schema $\Sigma^\ast$, we use In-Context Learning with elaborate prompts that exemplify the tree-based reasoning process. The prompts include:
\begin{itemize}[leftmargin=*]
    \item The complete operator space definition.
    \item Two few-shot examples demonstrating the \action{plan}--\action{expand}--\action{execute}--\action{answer} interaction protocol.
    \item The current source tables and target schema specification.
\end{itemize}

\stitle{Input Perturbation $\Phi$.}
To enhance robustness, we apply input perturbations $\Phi$ to the source tables before generating trajectories. Specifically, we:
\begin{itemize}[leftmargin=*]
    \item Randomly shuffle the order of columns in each source table.
    \item Randomly shuffle the order of rows in each source table.
    \item Randomly shuffle the order of source tables when multiple tables are present.
\end{itemize}
These perturbations ensure that the model does not rely on spurious positional correlations and learns to identify the correct columns and tables based on semantic content rather than position.

\stitle{Process Reward Filtering.}
The generated candidate trajectories are filtered using the process reward function $R_{\text{llm}}(\tau)$ (Section~\ref{subsec:rl_reward}). We retain only trajectories with $R_{\text{llm}}(\tau) > 0.8$ to ensure high training quality. This filtering process reduces the initial ~20,000 candidate trajectories to 11,198 high-quality samples.

\subsection{Environment Token Masking}
\label{subsec:appendix_token_masking}

As described in Section~\ref{subsec:rl_reward}, each trajectory contains both agent-generated tokens and environment-generated tokens. We provide additional details on the binary token mask implementation.

\stitle{Masking Mechanism.}
In our interaction protocol, a trajectory $\tau = (r_1, \dots, r_m)$ interleaves agent actions and environment observations. Specifically:
\begin{itemize}[leftmargin=*]
    \item \textbf{Agent-generated tokens} (unmasked, $M_t = 1$): Tokens within \action{plan}, \action{expand}, and \action{answer} blocks. These correspond to the agent's decision-making outputs and should be optimized by the policy gradient.
    \item \textbf{Environment-generated tokens} (masked, $M_t = 0$): Tokens within \action{execute} blocks. These are deterministic execution outputs (\eg serialized intermediate tables, error traces) that are not under the agent's control and should not contribute to policy gradient updates.
\end{itemize}

The binary mask $M_t$ is applied to the loss function in Eq.~(\ref{eq:sft_reasoning_learn}) during SFT and to the policy gradient objective during RL. This ensures that the model is only optimized for its own decision-making, not for predicting environment outputs.

\subsection{Hybrid Reward Details}
\label{subsec:appendix_hybrid_reward}

We provide additional details on the hybrid reward design described in Section~\ref{subsec:rl_reward}.

\stitle{Hyperparameters.}
The hybrid reward weights are set to $\alpha = 1$, $\beta = 1$, and $\gamma = 1$, \ie the total reward is the direct sum of the three components: $R(\tau) = R_{\text{out}}(\tau) + R_{\text{part}}(\tau) + R_{\text{llm}}(\tau)$. Note that $R_{\text{out}}$ takes a value of 2.0 for exact match and 0.0 otherwise, while $R_{\text{part}} \in [0, 1]$ and $R_{\text{llm}} \in [0, 1]$. This design gives the outcome reward a dominant weight (up to 2.0) to strongly incentivize exact correctness, while the partial and process rewards provide denser learning signals for trajectories that do not achieve exact match.

\stitle{LLM-as-Judge Scoring.}
The process reward $R_{\text{llm}}(\tau)$ is computed by an instruction-tuned LLM (\eg GPT-4o) that evaluates each trajectory on three criteria:
\begin{itemize}[leftmargin=*]
    \item \textbf{Plan--action consistency} (weight: 1/3): Whether the generated operator pipeline (from \action{expand}) correctly implements the plan (from \action{plan}).
    \item \textbf{Feedback responsiveness} (weight: 1/3): Whether subsequent \action{plan} outputs explicitly respond to the prior execution feedback (\eg reported error traces).
    \item \textbf{Backtracking justification} (weight: 1/3): Whether parent-node switches are supported by recorded failure evidence from the current branch.
\end{itemize}
Each criterion is scored on a scale of 0 to 1, and the final $R_{\text{llm}}(\tau)$ is the average of the three scores. The LLM judge is prompted with the full trajectory including all \action{plan}, \action{expand}, and \action{execute} blocks, along with the evaluation rubric.
